\section{Implementação}
\label{sec:implementacao}

Esta seção registra o estado observado no repositório em 29 de julho de 2026,
sem apresentar componentes inexistentes como concluídos. O repositório está
dividido em artigo, pesquisa e protótipo. No protótipo, o backend concentra
regras e casos de uso; a PWA consome a API ou executa um modo demonstrativo
sintético; e integrações externas respeitam as interfaces
\path{VehicleDataProvider} e \path{DispenserGateway}.

O backend FastAPI já recebe casos sintéticos, normaliza combustíveis, avalia a
compatibilidade, registra decisões e eventos minimizados em SQLite e aciona um
dispensador em memória. Antes de \path{UNLOCK}, a intenção
\path{PENDING_DEVICE} e o identificador do comando são persistidos; o ACK
síncrono do gateway simulado promove o estado para \path{AUTHORIZED}. Na
conclusão ou expiração, o backend tenta \path{LOCK_ALL}; se essa confirmação
falhar, o TTL local limita a liberação residual. O override é limitado a motivos
específicos de \path{REVIEW_REQUIRED}; \path{DENIED}, falha do dispositivo e
estado expirado não podem ser reativados. A proteção atual usa token Bearer
compartilhado, não autenticação individual.

O backend persiste \path{client_request_id} e um fingerprint HMAC da solicitação
para que uma reconciliação não provoque novo acionamento. Também mantém, somente
em memória e por TTL, um cache indexado por HMAC da placa para reutilizar a
consulta de compatibilidade no percurso PWA--autorização. O segredo HMAC é
configuração externa ao repositório. O MVP aceita, por padrão, apenas origens do
próprio computador; essa restrição e o token global não substituem autenticação
individual, RBAC, TLS ou gestão de sessão.

A PWA em React implementa os fluxos de placa, seleção, decisão, contagem de TTL,
conclusão, revisão de supervisor e reconciliação manual após resposta perdida,
sem repetir automaticamente comandos cujo efeito físico seja desconhecido. O
firmware experimental contém contrato MQTT,
cinco saídas distintas, sessão de inicialização, sequência monotônica persistida,
TTL local, rejeição de duplicatas e lógica desenergizada-igual-a-bloqueada.
Entretanto, o backend ainda utiliza apenas o gateway em memória: não há
integração MQTT ponta a ponta, placa ESP32 programada ou atuador físico ensaiado.
O modo SENATRAN é deliberadamente recusado na inicialização e não realiza
chamada externa.

\begingroup
\renewcommand{\arraystretch}{1.04}
\begin{table}[H]
  \centering
  \caption{Estado dos incrementos de implementação}
  \label{tab:implementacao-status}
  \footnotesize
  \begin{tabularx}{\textwidth}{@{}L{2.8cm}YL{3.2cm}@{}}
    \toprule
    \textbf{Incremento} & \textbf{Critério de conclusão} & \textbf{Estado nesta versão} \\
    \midrule
    Domínio e simuladores & Decisões reproduzíveis e bloqueio após término, expiração ou falha & Implementado no backend simulado \\
    API e persistência & Contratos, idempotência, auditoria minimizada e testes automatizados & FastAPI e SQLite implementados \\
    PWA e autenticação & Fluxos de frentista e supervisor utilizáveis em tablet & PWA implementada; autenticação individual pendente \\
    MQTT e bancada & Comandos idempotentes, TTL local e ciclos físicos repetidos & Firmware parcial; integração e bancada pendentes \\
    SENATRAN homologada & Caso de uso aprovado, contratos e certificados fora do repositório & Autorização pendente \\
    \bottomrule
  \end{tabularx}
  \\[2pt]\fontepropria
\end{table}
\endgroup

Na verificação local desta versão, \texttt{pytest} executou 82 testes do backend,
todos aprovados, com 93\% de cobertura de instruções. Na PWA, os 12 testes
automatizados foram aprovados e a compilação de produção terminou sem erro. A
lógica C++ portátil de intertravamento também compilou com avisos tratados como
erro e concluiu seu conjunto de asserções. A compilação PlatformIO para o alvo
\path{esp32dev} terminou sem erro; a gravação e execução em uma placa, a
integração entre componentes e os ciclos físicos não foram executados. Esses
resultados verificam incrementos de software, mas não equivalem à matriz integral
de aceitação de H1.

\pendente{concluir autenticação individual e RBAC, implementar o gateway MQTT no
backend, gravar e ensaiar o firmware em hardware aprovado pelo orientador e
obter autorização externa antes da integração SENATRAN.}

As tecnologias efetivamente usadas nesta etapa são FastAPI/Python, React/PWA,
SQLite e C++ para ESP32. MQTT está implementado apenas no lado experimental do
firmware. PostgreSQL e Docker permanecem possibilidades futuras. Qualquer ajuste
deverá preservar interfaces, estados seguros e testes de aceitação.
